%% ma: L12-B04-0018
%% lop: 12
%% chuong: 1
%% bai: 4
%% dang: 12.4.15
%% loai: DS
%% muc: [TH, NB, TH, TH]
%% dap_an: "SSĐS"
%% nguon: "HK1-DeOn_1_1.pdf (M10, Đề ôn tập học kì 1 Toán 12, Đề số 2), Phần II Câu 4"
%% kien_thuc: "Bài 1 (tính đơn điệu và cực trị), Bài 2 (giá trị lớn nhất, nhỏ nhất), Bài 3 (đường tiệm cận), Bài 4 (khảo sát và vẽ đồ thị hàm số)"
%% trang_thai: cho-duyet
%% ly_do: "Dạng mới đề xuất 12.4.15"
\begin{ex}
Cho hàm số $y=f(x)=\dfrac{-2x-3}{x+3}$.
\choiceTF{Hàm số đồng biến trên các khoảng $(-\infty;-3)$ và $(-3;+\infty)$}{Đồ thị hàm số có tiệm cận ngang $y=-3$}{\True Khoảng cách từ tâm đối xứng của đồ thị hàm số đến trục hoành bé hơn $3$}{Giá trị nhỏ nhất của hàm số đã cho trên đoạn $[0;2025]$ là $f(0)$}
\loigiai{a) Sai: $f'(x)=\dfrac{-3}{(x+3)^2}<0$, hàm số nghịch biến trên từng khoảng xác định.

b) Sai: tiệm cận ngang là $y=-2$.

c) Đúng: tâm đối xứng $I(-3;-2)$ cách trục hoành $2<3$.

d) Sai: hàm số nghịch biến trên $[0;2025]$ nên giá trị nhỏ nhất là $f(2025)$.}
\end{ex}
